\section{German Influence on Indology}

\begin{quotex}
This is the final chapter of \textbf{Rene Guenon}`s \emph{Introduction to the Study of Hindu Doctrines}, which was first published in 1921, right after the end of World War I. It was removed, at the author's request, from subsequent editions, including the version translated into English. It is quite harsh on German scholarship and even on Germans themselves. Moreover, his criticism extends to Europeans in general whose ``imagined superiority leads to misbehavior in the world."

I don't believe this chapter is available in English. Even if it is, it is worth repeating. 

\end{quotex}
\paragraph{German Influence}
It is rather curious to notice that the first Indologists, who were mostly English, without showing any very deep understanding, often said more correct things than those who came after them; no doubt they made many mistakes as well, but which at least were not of a systematic character, and which did not result from a bias, not even unconscious. The English mentality, of course, has no aptitude for metaphysical conceptions, but it does not make any pretension in this respect either, while the German mentality, which is not really better endowed, has the greatest illusions. To realize this, we need only compare what the two peoples have produced in terms of philosophy.

The English mind hardly left the practical order, represented by morality, sociology, and experimental science, represented by the science of psychology which it invented. When the English mind is concerned with logic, it is above all induction that he has in view and to which he gives preponderance over deduction. On the other hand, if we consider German philosophy, we only find in it hypotheses and systems with metaphysical pretensions, deductions from a fanciful starting point, ideas which might seem to be profound when they are simply nebulous; and this pseudo-metaphysics, which is everything that is farthest from true metaphysics, the Germans claim to find in others, whose conceptions they always interpret according to their own.

Nowhere is this latter fault more invincible than with them, because no other people has such a narrowly systematic mindset. Moreover, the Germans, in this regard, just push defects which are common to the whole European race to the extreme. Their national pride leads them to behave in Europe as Europeans in general, infatuated with their imaginary superiority, behave throughout the world. The extravagance is the same in both cases, with a simple difference in degree. It is therefore natural that the Germans imagine that their philosophers have thought all that it is possible for men to conceive, and no doubt they believe they are doing a great honor to other peoples by assimilating the conceptions of these to this philosophy of which they are so proud.

This does not prevent Schopenhauer from making a ridiculous disguise of Buddhism by making it a kind of ``pessimistic" moralism, and from giving the precise measure of his intellectual level by seeking ``consolations" in the ``Vedanta"; and we see, on the other hand, contemporary orientalists like Deussen claiming to teach the Hindus the true doctrine of Shankaracharya, to whom they quite simply lend the ideas of Schopenhauer! It is because the German mentality, by the very fact that it is an excessive form of the Western mentality, is the opposite of the East and cannot understand anything about it. As it nevertheless claims to understand, it necessarily distorts it. Hence these false assimilations against which we protest on all occasions, and in particular this application of the labels of modern Western philosophy to Eastern doctrines.

When one is incapable of doing metaphysics, it is certainly better not to be concerned with it, and positivism, in spite of its narrowness and incompleteness, seems to us still much preferable to the rantings of pseudo-metaphysics. The greatest mistake of the German orientalists is therefore to be unaware of their incomprehension, and to carry out works of interpretation which have no value, but which are imposed on all of Europe and very easily succeed in doing making themselves authoritative, because other peoples have nothing to oppose or to compare to it, and also because these works surround themselves with an apparatus of erudition which strongly impresses people who have a strong superstitious respect to for certain methods.

These methods, moreover, are also of Germanic origin, and it would be quite unfair not to recognize in the Germans the very real qualities which they possess in the matter of erudition: the truth is that they excel in the composition of dictionaries, grammars, and those voluminous works of compilation and bibliography which require nothing more than memory and patience. It is extremely unfortunate that they did not fully specialize in this kind of work which is very useful to consult on occasion, and which, appreciably, saves time for those who are capable of doing other things. What is hardly less regrettable is that these same methods, instead of remaining the prerogative of the Germans, to whose temperament they were particularly suited, have spread to all European universities, and especially in France, where they are considered to be only ``scientific", as if science and erudition were one and the same thing. In fact, as a consequence of this deplorable state of mind, scholarship comes to usurp the place of true science.

The abuse of cultivated scholarship for itself, the false belief that it can suffice to impart understanding of ideas, all of this, among Germans, can still be understood and excused to some extent; but, among peoples who do not have the same special aptitudes, we can only see in them the effect of a servile tendency to imitation, a sign of an intellectual decadence which it is high time to remedy, if we do not want to let it turn into a definitive decline.

The Germans very skillfully prepared the intellectual supremacy they dreamed of, by imposing both their philosophy and their methods of scholarship. Their Orientalism is, as we have just said, a product of the combination of these two elements. What is remarkable is how these things have become instruments in the service of national ambition; it would be quite instructive, in this regard, to study how the Germans were able to take advantage of the fanciful hypothesis of ``Aryanism", which they had moreover not even invented.

We do not believe, for our part, in the existence of an ``Indo-European" race, even if we do not wish to persist in calling it ``Aryan", which makes no sense. But what is significant is that German scholars have given this supposed race the denomination of ``Indo-Germanic", and that they have taken all their care to make this hypothesis plausible by supporting it with multiple ethnological and above all philological arguments. We do not wish to enter into this discussion at this point. We will only point out that the real resemblance which exists between the languages \hspace{0pt}\hspace{0pt}of India and Persia and those of Europe is by no means the proof of a community of race.

To explain it, it suffices that the ancient civilizations that we know of were originally brought to Europe by some elements relating to the source from which the Hindu and Persian civilizations directly proceeded. We know, in fact, how easy it is for a tiny minority, under certain conditions, to impose its language, with its institutions, on the mass of a foreign people, even though it is ethnically absorbed in it in a short time. A striking example is that of the establishment of the Latin language in Gaul, where the Roman presence, except in a few southern regions, was never more than negligible. The French language is undoubtedly of more or less pure Latin origin, and yet the Latin elements entered only a very small part in the ethnic formation of the French nation; the same is also true for Spain. On the other hand, the ``Indo-Germanic" hypothesis is all the less valid since the Germanic languages \hspace{0pt}\hspace{0pt}have no more affinity with Sanskrit than other European languages. It can only serve to justify the assimilation of Hindu doctrines to German philosophy; but, unfortunately, this supposition of an imaginary kinship does not stand the test of the facts, and nothing is in reality more dissimilar than a German and a Hindu, intellectually as well as physically, if not even in more ways.

The conclusion that emerges from all of this is that, in order to obtain interesting results, it is necessary to first get rid of this influence which for too long has weighed so heavily on Orientalism; and, although it is hardly possible for some individuals to break free from methods which constitute for them inveterate mental habits, we hope that, in general, recent events will be a favorable occasion for this liberation. However, let us understand our thought: if we want the disappearance of German influence in the intellectual field, it is because we consider it harmful in itself, and independent of certain historical contingencies which do not exist and change nothing. So it is not these contingencies that make us wish that the influence in question disappear, but we must take advantage of the state of mind that they have determined.

In the intellectual order, the only one we are dealing with here, sentimental concerns do not have to be taken into consideration. German conceptions today are worth exactly what they were worth a few years ago, and it is ridiculous to see men who had always professed a boundless admiration for German philosophy to suddenly start to denigrate it under the pretext of a patriotism which has nothing to do with these things. Fundamentally, this is little better than more or less consciously altering scientific or historical truth for reasons of national interest, precisely as the Germans are accused of doing.

For us, who owe nothing to Germanic intellectuality, who have never had the slightest esteem for pseudo-metaphysics in which it delights, and who have never accorded to erudition and its special procedures more than the most relative value and importance, we are very much at our ease to say what we think about it. And we would have said absolutely the same thing even if the circumstances had been quite different, but perhaps with less chance of agreeing in this with a generally prevalent trend. We will only add that, especially as regards France, what is at present most to be feared is that one escapes German influence only to fall under other influences which would be hardly less fatal. Reacting against the spirit of imitation therefore appears to us to be one of the first conditions for a true intellectual recovery: it is not a sufficient condition, no doubt, but it is at least a necessary, and even indispensable, condition.



\flrightit{Posted on 2021-07-04 by Cologero }

\begin{center}* * *\end{center}

\begin{footnotesize}\begin{sffamily}



\texttt{Arthur Konrad on 2021-07-05 at 05:44 said: }

First things first:

Whosoever aspires to be dabbling in the black arts must at the very least be able to read a language in which it is possible to communicate the music of things. Which is to say, Germans, as well as the Slavic peoples (with whom they share many mental attributes) are at a disadvantage from the onset. Among the latter, only the Russian language is an exception, and how perceptive of otherworldly things have the Russians always been! The rest of them have been derived from the speech of the folk, apparently, to satisfy the needs of both practicality and national-romantic fervour.

But alas, the `language of the people' is for issuing curt agricultural commands and managing domestic affairs, not for communicating the highest mystical teaching. That is why among the Eastern Slavs, the clergy thought it prudent to preserve the Church Slavonic, a true liturgical language. People like to say `Oh, but the German poetry!'. But German poetry is valued for its intellectual content, which is only one part of poetry. That being said, German is the most bastardized of all Germanic tongues, which really goes against all common wisdom, and conversely, the Russian is the most pristine of the Slavic ones.

As for me, if I haven't learned English my mental development would have been altogether stunted, what with my ears condemned to hearing only that which can be communicated in this clumsy, unwieldy linguistic apparatus which is called my mother tongue. For that I can thank globalism and capitalism.


\hfill

\texttt{Max on 2021-07-05 at 18:05 said: }

Perhaps it could be profitable to narrowly systematize Guénon's pure Latin thoughts, if not even fitting them into some form of compilation. Takeaways:

The German mentality has the greatest illusions.

The Germans always interpret other's conceptions according to their own.

No other people has such a narrowly systematic mindset.

The Germans imagine that they have thought everything that it is possible for men to conceive.

The Germans believe that they are doing a great honour by assimilating other people's conceptions.

The other more superstitious peoples of Europe are easily fooled by German methods of erudition.

It is unfortunate that the Germans did not specialize in being Guénon's personal secretary – they could have been useful on occasion.

Germans excel in things that require nothing more than memory and patience.

Scholarship has come to usurp true science.

The Germans very skillfully prepared the intellectual supremacy of their dreams.

Guénon knows how easy it is for a tiny minority to impose its language and institutions on a mass of a foreign people.

Germanic has no more affinity with Sanskrit than other European languages.

The kinship is an imaginary supposition anyway.

Intellectually, nothing is more dissimilar than a German and a Hindu.

Physically, nothing is more dissimilar than a German and a Hindu.

If not even in more ways, nothing is more dissimilar than a German and a Hindu.

It is hardly possible for some individuals to break free from Germanic methods, which for them constitute inveterate mental habits.

Guénon hopes that certain events will nevertheless present a favorable occasion for the liberation from the German mentality.

German influence is harmful independently of certain historical contingencies, which do not exist, and would change nothing even if they did.

Guénon owes nothing to Germanic intellectuality.

In the intellectual order, sentimental concerns do not have to be taken into consideration.


\end{sffamily}\end{footnotesize}
